\section{Zero-Gap Nine-Point Enclosure}
\label{sec:strategy4-reader}
\label{sec:zero-gap-nine-point-optimization}
\label{sec:ab-core}

We first compute the local frontier, then construct the forced witnesses.
The enclosure argument reduces to four supporting contacts, including two
exact disk-tangent certificate bounds.

The remaining zero-gap case has one supercritical V triangle. There is no
boundary gap to use as a witness, so the obstruction must be built inside
$H$. The construction and its verification have three stages.

\readerheading{Force the points.}
Strict boundary handoffs give common lower bounds for the boundary reaches.
These force six radial points into the C triangle. Three further points are
chosen on the exact frontier of the feasible union of the supercritical
triangle; separate diameter and handoff arguments exclude every other V
triangle. Exclusion from the supercritical triangle alone would not be enough.
A frontier point may belong to a closed feasible triangle while being missed
by its open interior.

\readerheading{Simplify the set to be enclosed.}
The convex hull of the six radial points contains a disk. Replacing those
six points by this disk gives a smaller comparison set. In the nontrivial
range, the two outer frontier points are then moved inward by a single
Newton step at the fixed line parameter $k=1/2$. Convexity keeps these comparison points inside any candidate
containing the original witnesses. These are exact inner comparisons,
not numerical approximations to a covering configuration.

\readerheading{Bound the enclosing triangle.}
The inner comparison reduces the enclosure problem to four supporting
contacts: two contacts along consecutive-point lines and two disk-tangent
contacts. The line estimates are analytic. The remaining paired tangent
inequalities are reduced to the exact polynomial certificate. At the end,
the comparison set requires side at least one, contradicting its containment
in the open C triangle. The calculations below justify each of these stages;
the certificate enters only at the last one.

\readerheading{Active tangent calculation.}
The two outer comparison points use the fixed-start formulas, not the older
junction-start formulas. The exact coefficient comparisons and radius envelope
are specified in \zcref{thm:paper-exact-mixed-certificate}.
The actual disk is unchanged; the smaller scalar radius is used only for
initially checking the tangent residuals.


\subsection{Corner-clipped $AB$-unions}

In this subsection the corner is the local origin $O_{\rm loc}$, not the
hexagon center $O$. The local anchors $A,B$ used here are also distinct
from the Newton comparison points introduced later. The corner chart
\[
 \Psi_4(u,v)=V_4+u(V_5-V_4)+v(V_3-V_4)
\]
has coordinate axes corresponding to the forward and backward edge
directions. We identify a local vector $ue_A+ve_B$ with its pair $(u,v)$
when applying this chart.

Let $e_A,e_B$ be unit vectors with angle $120^\circ$, put
\[
 O_{\rm loc}=0,\qquad A=ae_A,\qquad B=be_B,
 \qquad
 \mathcal C_{\rm loc}=\{ue_A+ve_B:u,v\ge0\},
\]
and define
\begin{equation}
 \mathcal R_{AB}^{\rm loc}(a,b)=
 \mathcal C_{\rm loc}\cap\bigcup\left\{T:
 \begin{array}{l}
 T\text{ is a closed unit equilateral triangle},\\
 O_{\rm loc},A,B\in T
 \end{array}\right\}.
 \label{eq:ab-union}
\end{equation}
The chart sends the forward and backward parameters to the hexagon.
In the covering application below, their local order is $(b,a)$.

The strict supercritical branch admits a simpler closed form.  In the next
statement write $x=ue_A+ve_B$ and
$\lVert x\rVert^2=u^2+v^2-uv$.

\begin{theorem}[Strict supercritical $AB$-union]
\label{thm:strict-ab-union}
Assume $a,b>0$ and
\[
 a+b>1,\qquad \rho=a^2+ab+b^2<1,
\]
put $h=\sqrt3/2$, $D=\sqrt{4\rho-3}$, and define
\begin{align*}
 \alpha&=\frac{h(a+2b-aD)}{2\rho},&
 \beta&=\frac{h(a-b+(a+b)D)}{2\rho},\\
 \gamma&=\frac{h(-a+b+(a+b)D)}{2\rho},&
 \delta&=\frac{h(2a+b-bD)}{2\rho}.
\end{align*}
Then all four coefficients are positive and
\begin{equation}
 \mathcal R_{AB}^{\rm loc}(a,b)=
 (\mathcal C_{\rm loc}\cap\mathcal D_A\cap\mathcal H_A)
 \cup(\mathcal C_{\rm loc}\cap\mathcal D_B\cap\mathcal H_B),       \label{eq:strict-ab-identity}
\end{equation}
where
\begin{align*}
 \mathcal D_A&=\{(u,v):(u-a)^2+v^2-(u-a)v\le1\},\\
 \mathcal D_B&=\{(u,v):u^2+(v-b)^2-u(v-b)\le1\},\\
 \mathcal H_A&=\{(u,v):\alpha(u-a)+\beta v\le0\},\\
 \mathcal H_B&=\{(u,v):\gamma u+\delta(v-b)\le0\}.
\end{align*}
Its non-axis frontier consists, in counterclockwise order, of the
$A$-circle, the $A$-line, the $B$-line, and the $B$-circle.
\end{theorem}

\begin{proof}
The identities
\[
 4\rho-3(a+b)^2=(a-b)^2,
 \qquad
 (a+b)^2D^2-(a-b)^2=4\rho((a+b)^2-1)
\]
give $0<D<1$ and positivity of the four coefficients.

We apply the caliper criterion of \zcref{thm:cert-caliper}.  For a
point $x=(u,v)$ in the relative interior of $\mathcal C_{\rm loc}$ but outside
$\triangle O_{\rm loc}AB$, the four hull vertices occur in the order $O_{\rm loc},B,x,A$.
The two cone-axis hull edges cannot certify a unit triangle: their support
sums, divided by $h$, are respectively
\[
 \max\{a,u\}+\max\{b,v-u\},\qquad
 \max\{b,v\}+\max\{a,u-v\},
\]
and are at least $a+b>1$.  It remains to test the $Ax$ and $Bx$ calipers.

For the $Ax$ caliper put $p=u-a$ and
$r^2=p^2+v^2-pv$.  Exact evaluation of the three supports gives
\begin{equation}
 \frac{G_A}{h}=
 \frac{\max\{r^2,-ap+(a+b)v\}+N_A}{r},
 \quad
 N_A=\begin{cases}
 0,&p\le0,\\
 ap,&0\le p\le v,\\
 (a+b)p-bv,&p\ge v.
 \end{cases}                                      \label{eq:Ax-caliper}
\end{equation}
No support label is omitted in these three sectors.  If $p\le0$,
$G_A\le h$ is equivalent to
\[
 r\le1,\qquad -ap+(a+b)v\le r.
\]
The second inequality is the half-plane condition in
\eqref{eq:strict-ab-identity}, because
\[
 r^2-(-ap+(a+b)v)^2
 =\frac{(cp+dv)(c'p+d'v)}{4\rho},
\]
where
\[
 c=a+2b-aD,\quad d=a-b+(a+b)D,
 \quad c'=a+2b+aD,\quad d'=a-b-(a+b)D,
\]
and $c,c',d>0>d'$.  Thus the low sector is exactly
$\mathcal D_A\cap\mathcal H_A$.  In the middle sector $0<p\le v$ one has
$r\le v$, whereas the numerator in \eqref{eq:Ax-caliper} is at least
$(a+b)v>r$, so no fit is possible.

In the high sector $p\ge v$, the inequalities $G_A\le h$ imply
\begin{equation}
 r^2+(a+b)p-bv\le r,\qquad bp+av\le r.          \label{eq:Ax-high}
\end{equation}
Write $x-A=ry(\theta)$, $0\le\theta\le60^\circ$, with cone coordinates
\[
 y(\theta)=te_A+we_B,\quad
 t=\frac2{\sqrt3}\cos(\theta-30^\circ),\quad
 w=\frac2{\sqrt3}\sin\theta.
\]
Then \eqref{eq:Ax-high} gives
\[
 r\le f(\theta)=1-(a+b)t+bw,\qquad S(\theta)=bt+aw\le1.
\]
Let $n_B$ be the unit vector with dual coordinates $(\gamma,\delta)$.
The identities
\[
 b\gamma+(a+b)\delta=h,\quad
 (a+b)\gamma+a\delta=\frac h2(1+D),\quad
 b\delta-a\gamma=\frac h2(1-D)
\]
show that the unique angle $\theta_B$ of this normal satisfies
$S(\theta_B)=1$ and $f(\theta_B)=(1-D)/2$.  The function $S$ has at most
one critical point, a maximum, so $S(\theta)\le1$ implies
$\theta\le\theta_B$.  On the feasible part $f,f'\ge0$; hence
$f(\theta)\cos(\theta_B-\theta)$ is increasing.  Consequently
\[
 \langle n_B,x-B\rangle
 \le a\gamma-b\delta+\frac h2(1-D)=0.
\]
The diameter bound also gives $\lVert x-B\rVert\le1$.  Thus every high
sector fit belongs to $\mathcal D_B\cap\mathcal H_B$.  Reflection exchanges
$a,u$ with $b,v$ and supplies the corresponding result for the $Bx$
caliper.  This proves \eqref{eq:strict-ab-identity} in the cone interior;
the axes follow by taking limits, since both sides are closed.

For completeness, the line--circle junctions are obtained by solving the
two displayed line and circle equations.  The line normals satisfy
\[
 \alpha^2+\alpha\beta+\beta^2=h^2,\qquad
 \gamma^2+\gamma\delta+\delta^2=h^2,
\]
and the line determinant is
\[
 \alpha\delta-\gamma\beta
 =\frac{3(1-D)(D+3)}{16\rho}>0.
\]
The successive cone slopes are
\[
 +\infty,\quad \frac2{1-D},\quad
 \frac{a(1-D)+2b}{2a+b(1-D)},\quad
 \frac{1-D}{2},\quad0;
\]
the two middle comparisons reduce to
$a(4-(1-D)^2)>0$ and $b(4-(1-D)^2)>0$.  This proves the asserted four-piece
order.
\end{proof}

\begin{figure}[!htbp]
\centering
\begin{minipage}[t]{.48\linewidth}
\centering
\resizebox{.96\linewidth}{!}{%
\begin{tikzpicture}[x=1.55cm,y=1.55cm,>=Latex]
  \node[panellabel,anchor=west] at (-1.05,2.18)
    {(a) exact strict-supercritical $AB$ union};
  \coordinate (O) at (0,0);
  \coordinate (A) at (2.25,0);
  \coordinate (B) at (-1.125,1.949);
  \coordinate (P1) at (1.78,.33);
  \coordinate (Qm) at (1.35,.66);
  \coordinate (Q0) at (.57,1.10);
  \coordinate (Qp) at (-.08,1.46);
  \coordinate (P4) at (-.62,1.78);
  \draw[black!45,line width=.75pt] (O)--(2.55,0);
  \draw[black!45,line width=.75pt] (O)--(-1.28,2.217);

  \fill[VertexOrange!13]
    (O)--(A).. controls (2.08,.10) and (1.93,.23)..(P1)--(Qm)--(Q0)--cycle;
  \fill[CenterBlue!12]
    (O)--(Q0)--(Qp)--(P4).. controls (-.82,1.88) and (-1.00,1.92)..(B)--cycle;

  \draw[DemandRed,line width=1pt]
    (A).. controls (2.08,.10) and (1.93,.23)..(P1);
  \draw[VertexOrange,line width=1pt] (P1)--(Q0);
  \draw[CenterBlue,line width=1pt] (Q0)--(P4);
  \draw[GapPurple,line width=1pt]
    (P4).. controls (-.82,1.88) and (-1.00,1.92)..(B);

  \fill (A) circle (.018) node[below=2pt,figlabel] {$A=ae_A$};
  \fill (B) circle (.018) node[left=2pt,figlabel] {$B=be_B$};
  \fill (O) circle (.018) node[below left=1pt,figlabel] {$O$};
  \node[figlabel,text=DemandRed,rotate=28] at (2.06,.25) {$A$-circle};
  \node[figlabel,text=VertexOrange,rotate=30] at (1.30,.72) {$A$-line};
  \node[figlabel,text=CenterBlue,rotate=31] at (.12,1.44) {$B$-line};
  \node[figlabel,text=GapPurple,rotate=21] at (-.78,1.93) {$B$-circle};
  \node[figlabel,align=center] at (.45,-.55)
    {$\mathcal R(a,b)=(C\cap\mathcal D_A\cap\mathcal H_A)
      \cup(C\cap\mathcal D_B\cap\mathcal H_B)$};
\end{tikzpicture}%
}
\end{minipage}\hfill
\begin{minipage}[t]{.48\linewidth}
\centering
\resizebox{.96\linewidth}{!}{%
\begin{tikzpicture}[x=1.55cm,y=1.55cm,>=Latex]
  \node[panellabel,anchor=west] at (-1.05,2.18)
    {(b) the three asymmetric frontier witnesses};
  \coordinate (O) at (0,0);
  \coordinate (A) at (2.25,0);
  \coordinate (B) at (-1.125,1.949);
  \coordinate (P1) at (1.78,.33);
  \coordinate (Qm) at (1.35,.66);
  \coordinate (Q0) at (.57,1.10);
  \coordinate (Qp) at (-.08,1.46);
  \coordinate (P4) at (-.62,1.78);
  \draw[black!45,line width=.75pt] (O)--(2.55,0);
  \draw[black!45,line width=.75pt] (O)--(-1.28,2.217);
  \draw[black!55,line width=.75pt]
    (A).. controls (2.08,.10) and (1.93,.23)..(P1)--(Q0)--(P4)
    .. controls (-.82,1.88) and (-1.00,1.92)..(B);
  \fill[DemandRed] (Qm) circle (.035) node[below right=2pt,figlabel] {$Q_-$};
  \fill[DemandRed] (Q0) circle (.035) node[right=2pt,figlabel] {$Q_0$};
  \fill[DemandRed] (Qp) circle (.035) node[above right=2pt,figlabel] {$Q_+$};

  \draw[black!35,dashed] (Qm) circle (.76);
  \draw[black!35,dashed] (Qp) circle (.76);
  \draw[flow] (Qm)--($(Qm)!1.42!(O)$);
  \draw[flow] (Q0)--($(Q0)!1.42!(O)$);
  \draw[flow] (Qp)--($(Qp)!1.42!(O)$);
  \node[figlabel,atlasbox,align=left,anchor=west] at (.80,1.57)
    {frontier excludes the actual supercritical role;\\
     moving unit-circle signs exclude the two adjacent roles;\\
     vertex distances exclude the three nonlocal roles};
  \node[figlabel,align=center] at (.45,-.55)
    {$Q_-,Q_0,Q_+\in\operatorname{int}(H)
      \setminus\bigcup_{i=0}^{5}U_i$};
\end{tikzpicture}%
}
\end{minipage}
\caption{The strict-supercritical $AB$ frontier and the three asymmetric
witnesses.  The source-conditioned union has, in counterclockwise order, an
$A$-circle arc, an $A$-line segment, a $B$-line segment, and a $B$-circle
arc.  The points $Q_-,Q_0,Q_+$ are selected on this exact frontier and are
then excluded from the six actual V roles by frontier, moving-circle, and
vertex-distance arguments.  The drawing is schematic; the labelled frontier
order and witness roles are exact.}
\label{fig:strict-ab-frontier-witnesses}
\end{figure}


\subsection{Exact-one handoffs and common reach bounds}

\begin{lemma}[Exact-one handoff chain]
\label{lem:ab-extreme-jump}
Assume the six V triangles cover $\partial H$ and exactly one actual V triangle is
supercritical.  Rotate its index to $4$.  The strict handoffs supplied by
\zcref{prop:strict-handoffs} satisfy
\begin{equation}
 \xi_3<\xi_2<\xi_1<\xi_0<\xi_5<\xi_4.
 \label{eq:extreme-chain}
\end{equation}
With
\begin{equation}
 a=1-\xi_3,\qquad b=\xi_4,\qquad p=1-b,\qquad q=1-a,
 \label{eq:reader-ab-parameters}
\end{equation}
one has
\begin{equation}
 0<a,b<1,\qquad a+b>1,\qquad a^2+ab+b^2<1,       \label{eq:ab-strict-domain}
\end{equation}
Moreover $a=a_4<A_4$ and $b=b_4<B_4$, and every selected V triangle obeys
\[
 a_i\ge p,\qquad b_i\ge q.
\]
\end{lemma}

\begin{proof}
Every selected V triangle other than $4$ is strictly nonsupercritical, so
$\xi_i<\xi_{i-1}$ for $i\ne4$, which is precisely
\eqref{eq:extreme-chain}.  In particular $\xi_3$ is the global minimum and
$\xi_4$ the global maximum.  The two anchors of V triangle $4$ lie in its open triangle;
their mutual squared distance is $a^2+ab+b^2$, strictly below the diameter
squared of the closed unit triangle.  This proves \eqref{eq:ab-strict-domain}.
Finally every $\xi_j\in[\xi_3,\xi_4]$, so
$1-\xi_{i-1}\ge1-\xi_4=p$ and $\xi_i\ge \xi_3=q$.
\end{proof}

Define the two boundary anchors
\[
 P_A:=V_4+a(V_3-V_4),
 \qquad
 P_B:=V_4+b(V_5-V_4),
\]
and the feasible union
\[
 \mathcal R_{AB}(a,b):=
 \bigcup\left\{T\cap H:
 \begin{array}{l}
 T\text{ is a closed unit equilateral triangle},\\
 V_4,P_A,P_B\in T
 \end{array}\right\}.
\]

By the wedge reduction, this union equals
$\Psi_4(\mathcal R_{AB}^{\rm loc}(b,a))$.

\subsection{Six directly forced radial points}

Put $h=\sqrt3/2$ and
\[
 s=p+q,\qquad m=\min(p,q),\qquad M=\max(p,q),
 \qquad \chi=s^4-s^2+pq.
\]
The strict domain \eqref{eq:ab-strict-domain} gives
\begin{equation}
 pq>(1-s)(2-s),\qquad m>1-s,\qquad s>1-m.         \label{eq:pq-domain}
\end{equation}
Let $c_*=c_{\max}(p,q)$ be the exact radial envelope of
\zcref{prop:exact-admissible-set}.  On the present subcritical
domain, equation \eqref{eq:radial-envelope} becomes
\begin{equation}
 c_*=\begin{cases}
 C_L(m),&\chi\le0,\\[1mm]
 \displaystyle\frac{2M}{1+\sqrt{4s^2-3}},&\chi>0,
 \end{cases}                                    \label{eq:core-cstar}
\end{equation}
where $C_L(m)$ is the unique root in $[\sqrt3/2,1]$ of
\[
 z^4-z^2+mz-m^2=0.
\]
At $\chi=0$ the two expressions agree and equal $s$.

\begin{lemma}[Direct radial forcing]
\label{lem:symmetric-core-witness}
One has
\begin{equation}
 0<c_*<1,\qquad c_*>1-m.                         \label{eq:cstar-bounds}
\end{equation}
Assume that $U_C,U_0,\ldots,U_5$ cover $H$, with no restriction on the
normalized V types. Then, for every $i$,
\[
 D_i=(1-c_*)V_i\in U_C.
\]
Consequently $U_C$ contains the centered closed disk
\begin{equation}
 \mathcal D_\eta=\{x:\lVert x\rVert\le\eta\},
 \qquad \eta=h(1-c_*).                           \label{eq:forced-disk}
\end{equation}
\end{lemma}

\begin{proof}
The complementary pair is positive and has $p+q<1$.
With $m=\min(p,q)$ and $c_*=c_{\max}(p,q)$, \zcref{lem:bc-capacity-tools} gives
$1-c_*\ge m(1-m)/2>0$, while
\zcref{lem:strict-own-ray-positive-slack} gives $1-c_*<m$.
Hence $0<c_*<1$ and $c_*>1-m$.
Every selected pair dominates $(p,q)$ by \zcref{lem:ab-extreme-jump}.
Apply \zcref{cor:new-uniform-common-pair-forcing} to obtain the unchanged
six radial points and their inscribed disk.
\end{proof}

We next define the three asymmetric witnesses.  At $V_4$ use local
coordinates
\[
 P=\Psi_4(u,v),\qquad
 \mathsf q(u,v)=u^2+v^2-uv.
\]
The forward-reach lower bound is $b$ and the backward-reach lower bound is $a$.  Accordingly,
set
\begin{align*}
 \alpha&=\frac{h(b+2a-bD)}{2\rho},&
 \beta&=\frac{h(b-a+(a+b)D)}{2\rho},\\
 \gamma&=\frac{h(-b+a+(a+b)D)}{2\rho},&
 \delta&=\frac{h(2b+a-aD)}{2\rho},
\end{align*}
where now $\rho=a^2+ab+b^2$ and $D=\sqrt{4\rho-3}$.  The frontier coefficients here are not the signed-center slacks.
Applying \zcref{thm:strict-ab-union} with the exchanged pair $(b,a)$,
the two local line pieces are
\[
 L_-(\lambda)=(b-\beta\lambda,\alpha\lambda),\qquad
 L_+(\mu)=(\delta\mu,a-\gamma\mu).
\]
They meet at $\lambda_* =\mu_*=8h\rho/(3(D+3))$ in
\begin{equation}
 J=\left(\frac{2b+a-aD}{D+3},
          \frac{b+2a-bD}{D+3}\right).           \label{eq:line-junction}
\end{equation}
Put
\[
 E_-=(1-b,2-b),\qquad E_+=(2-a,1-a).
\]
The actual adjacent handoffs have local centers
\[
 Z_5(t)=(1+t,t),\qquad 1-a\le t\le b,
 \quad\text{and}\quad
 Z_3(y)=(y,1+y),\qquad 1-b\le y\le a.
\]
The restrictions of the two unit-circle equations are
\begin{align}
 g_-(\lambda)&=\frac34\lambda^2
 -3(\alpha+b\beta)\lambda+3b^2-3b+2,\notag\\
 g_+(\mu)&=\frac34\mu^2
 -3(\delta+a\gamma)\mu+3a^2-3a+2.              \label{eq:witness-quadratics}
\end{align}
Let $\lambda_-$ and $\mu_+$ be their first roots; explicitly,
\begin{align*}
 \lambda_-&=2(\alpha+b\beta)
 -\frac23\sqrt{9(\alpha+b\beta)^2-9b^2+9b-6},\\
 \mu_+&=2(\delta+a\gamma)
 -\frac23\sqrt{9(\delta+a\gamma)^2-9a^2+9a-6}.
\end{align*}

\begin{lemma}[Normalized two-coefficient bounds]
\label{lem:F-normal-pair-bounds}
Suppose $x,y>0$, $x^2+xy+y^2=1$, $r>1$, $0<c<1$, and
$rx+cy=1$. Then
\[
 0<x,y<1,\qquad cx<\frac1{\sqrt3},\qquad x-y<2(1-c).
\]
\end{lemma}
\begin{proof}
Since $cy<1-x$ and $1-x^2=y(x+y)$,
\[
 c<\frac{1-x}{y}=\frac{x+y}{1+x}.
\]
Consequently
\[
 cx<\frac{x(x+y)}{1+x}\le\frac{x+y}{2}\le\frac1{\sqrt3}.
\]
The last inequality follows from
$1=x^2+xy+y^2\ge3(x+y)^2/4$.
Also $1-c>(1-y)/(1+x)$. If $x<y$, the second claimed
bound is immediate; otherwise
\[
 x-y\le1-y<(1+x)(1-c)<2(1-c).
\]
\end{proof}

\begin{lemma}[Moving adjacent-triangle signs]
\label{lem:fixed-line-signs}
Throughout \eqref{eq:ab-strict-domain},
\begin{equation}
 \lambda_*<1/\sqrt3<\lambda_-<h^{-1},\qquad
 \mu_*<1/\sqrt3<\mu_+<h^{-1}.
 \label{eq:first-root-order}
\end{equation}
Moreover $J_u,J_v<1/2$. The point $L_+(\mu_+)$ satisfies the stronger bound
\[
 (L_+(\mu_+))_u+(L_+(\mu_+))_v<2-a<3-2a,
\]
and the reflected bound holds for $L_-(\lambda_-)$. Writing
$F(P,t)=\mathsf q(P-Z_5(t))-1$, one has, for every $t\in[1-a,b]$,
\begin{equation}
 F(L_+(\mu_+),t)\ge0,\qquad F(J,t)>0,\qquad
 F(L_-(\lambda_-),t)>0.
 \label{eq:moving-right-signs}
\end{equation}
Under the reflection $(u,v,a,b)\leftrightarrow(v,u,b,a)$,
the analogous signs hold for every $y\in[1-b,a]$: the first-root point
$L_-(\lambda_-)$ has nonnegative sign against $Z_3(y)$, and the other
two witnesses have strictly positive sign.
\end{lemma}

\begin{proof}
Put $U=a+b-1$, $z=a-b$, and normalize the line coefficients by
\[
 x=\alpha/h,\quad y=\beta/h,\quad
 w=\delta/h,\quad v=\gamma/h.
\]
Direct substitution in their definitions gives
\[
 x^2+xy+y^2=w^2+wv+v^2=1,
 \qquad (a+b)x+by=(a+b)w+av=1.
\]
All four coefficients are positive. Also
$0<U<2/\sqrt3-1<1/6$ and $0<D<1$.
Lemma~\ref{lem:F-normal-pair-bounds} therefore gives
\begin{equation}
 bx<1/\sqrt3,\qquad aw<1/\sqrt3,\qquad
 x-y<2(1-b),\qquad w-v<2(1-a).
 \label{eq:F-normal-pair-consequences}
\end{equation}

Reparametrize the two frontier lines by $k=h\lambda$ and $k=h\mu$:
\[
 P_-(k)=(b-yk,xk),\qquad P_+(k)=(wk,a-vk).
\]
Their common junction parameter is
\[
 k_*:=h\lambda_*=h\mu_*=
 \frac{2\rho}{D+3}=\frac{D^2+3}{2(D+3)},
 \qquad \frac38<k_*<\frac12.
\]
The circle restrictions become
\[
 \begin{aligned}
 f_b(k)&:=\mathsf q(P_-(k)-E_-)-1
       =k^2-3(1-Ux)k+3b^2-3b+2,\\
 f_a(k)&:=\mathsf q(P_+(k)-E_+)-1
       =k^2-3(1-Uw)k+3a^2-3a+2.
 \end{aligned}
\]
For the first quadratic,
\[
 f_b(1/2)=3(b-1/2)^2+\tfrac32Ux>0,
 \qquad f_b'(1/2)=-2+3Ux<-\tfrac32.
\]
Moreover
\[
 f_b(1)=3\{Ux-b(1-b)\}<0,
\]
since $\rho=1+U(1+a)-b(1-b)<1$ and $x<1<1+a$.
Thus its first root $k_-=h\lambda_-$ lies in $(1/2,1)$.
The reflected calculation gives $k_+=h\mu_+\in(1/2,1)$.
Both quadratics decrease on $[k_*,1/2]$, so
\[
 f_b(k_*)>f_b(1/2)>0,\qquad f_a(k_*)>f_a(1/2)>0.
\]
This proves the root ordering and the positive junction values without
expanding a junction polynomial.

For the junction coordinates, use
\[
 D^2-(3U-z)^2=6U(1+z-U)>0,\qquad
 D^2-(3U+z)^2=6U(1-z-U)>0.
\]
The first identity implies $3U-z<D(1+2a)$, which is exactly
$J_u<1/2$; the second gives $J_v<1/2$. Hence $J_u+J_v<1$.
Because
\[
 \partial_tF(P,t)=1+2t-P_u-P_v,
\]
$F(J,t)$ is strictly increasing for $t\ge0$.
Its value at $t=1-a$ is $f_a(k_*)>0$, proving the junction sign.

On the opposite line, direct differentiation gives
\[
 \frac{d}{dk}F(P_-(k),t)
 =2k+(1-b)(2y+x)+t(y-x).
\]
For $k\ge k_*$ and $t\in[1-a,b]$, this is at least
\[
 2k-bx>\frac34-\frac1{\sqrt3}>0.
\]
Thus $F(P_-(k_-),t)>F(J,t)>0$ throughout the handoff interval.

On the correct line, $0<k_+<1$ and \eqref{eq:F-normal-pair-consequences}
give
\[
 (P_+(k_+))_u+(P_+(k_+))_v
 =a+k_+(w-v)<2-a.
\]
Indeed, if $w-v<0$ the sum is below $a<2-a$; otherwise it is below
$a+2(1-a)$. Therefore, for $t\ge1-a$,
\[
 \partial_tF(P_+(k_+),t)>1+2(1-a)-(2-a)=1-a>0.
\]
As $F(P_+(k_+),1-a)=f_a(k_+)=0$, its sign is nonnegative
throughout the interval and positive away from its left endpoint.
Reflection proves all the remaining assertions.
\end{proof}

Define the local and global witnesses by
\begin{equation}
 Q_-^{\rm loc}=L_-(\lambda_-),\qquad
 Q_0^{\rm loc}=J,\qquad
 Q_+^{\rm loc}=L_+(\mu_+),                       \label{eq:asym-witnesses}
\end{equation}
and put $Q_j=\Psi_4(Q_j^{\rm loc})$ for $j\in\{-,0,+\}$.

\begin{lemma}[Direct asymmetric forcing]
\label{lem:asymmetric-core-witness}
Assume that $U_C,U_0,\ldots,U_5$ cover $H$.  Then
\[
 Q_-,Q_0,Q_+
 \in\interior(H)\setminus\bigcup_{i=0}^5U_i,
\]
and consequently all three witnesses belong to $U_C$.
\end{lemma}

\begin{proof}
Put $\bar\alpha=\alpha/h$, $\bar\beta=\beta/h$, and $\bar\gamma=\gamma/h$.
Write $x=\bar\alpha$, $y=\bar\beta$. The coefficient identities give
$\rho-1=(y-b)(x+y-b)/x^2$. Since $x+y>1>b$, one has $y<b$;
reflection gives $\bar\gamma<a$. Thus both outer line endpoints at
$k=1$ have coordinates in $(0,1)^2$. The junction and the two first-root
points lie on the intervening open segments, so $0<u,v<1$ for all three points.  Each point has local coordinates in $\mathcal R_{AB}^{\rm loc}(b,a)$ and hence
lies in a closed unit triangle containing $V_4$, so
$\mathsf q(u,v)\le1$; hence
$|u-v|<1$ and all three points lie in $\interior(H)$.

In local coordinates the closure of $U_4$ is represented in
$\mathcal R_{AB}^{\rm loc}(b,a)$, since $U_4$ contains the actual handoffs
$X_3(\xi_3),X_4(\xi_4)$ and its distinguished vertex $V_4$.  All three witnesses lie
on its exact non-axis frontier.  If one belonged to the open triangle $U_4$,
a sufficiently small plane ball about it would lie both in $U_4$ and in the
interior of the corner cone, making the witness an interior point of
$\mathcal R_{AB}^{\rm loc}(b,a)$. This contradiction excludes $U_4$.

For the triangles $U_0,U_1,U_2$, their contained vertices have local coordinates
\[
 \begin{array}{c|ccc}
 i&0&1&2\\ \hline
 V_i&(2,1)&(2,2)&(1,2).
 \end{array}
\]
For $P=(u,v)$,
\[
 \lVert P-(2,1)\rVert^2-1=\mathsf q(u,v)+2-3u,
 \quad
 \lVert P-(1,2)\rVert^2-1=\mathsf q(u,v)+2-3v.
\]
The inequalities $J_u,J_v<1/2$ and the line order give distance greater than
one from $V_0$ for $Q_-,Q_0$ and from $V_2$ for $Q_0,Q_+$.  For the remaining
pair, the first-root circle and the sum bound in
\zcref{lem:fixed-line-signs} give
\[
 \lVert Q_+-V_0\rVert^2-1
 =a(3-a-(Q_+)_u-(Q_+)_v)>a^2,
\]
and, by reflection, $\lVert Q_--V_2\rVert^2-1>b^2$.  All three points are
farther than one from $V_1=(2,2)$, because writing $u=1-\xi$, $v=1-\eta$
gives
\[
 \lVert P-V_1\rVert^2
 =1+\xi+\eta+\xi^2+\eta^2-\xi\eta>1.
\]
Since two points in one open unit equilateral triangle have distance strictly
less than one, these inequalities exclude $U_0,U_1,U_2$.

The actual handoff $X_5(\xi_5)\in U_5$ has local coordinates
$Z_5(\xi_5)$, and
$1-a<\xi_5<b$ by \eqref{eq:extreme-chain}.  The three signs in
\eqref{eq:moving-right-signs} therefore put every witness at distance at
least one from $X_5(\xi_5)$, excluding $U_5$.  Similarly, if $y=1-\xi_2$, then
$1-b<y<a$ and $X_2(\xi_2)\in U_3$ has local coordinates $Z_3(y)$;
the reflected signs exclude $U_3$.
Thus no V triangle contains any witness.  The assumed cover now forces all
three interior witnesses into $U_C$.
\end{proof}

Set $r_*:=\eta=h(1-c_*)$ and $P_i^{\rm rad}:=(1-c_*)V_i$. Then
\begin{equation}
 \mathbb B(O,r_*)\subseteq\conv\{P_0^{\rm rad},\ldots,P_5^{\rm rad}\}
 \subset U_C.
 \label{eq:reader-forced-disk}
\end{equation}
The three additional witnesses satisfy
\begin{equation}
 Q_-,Q_0,Q_+\in U_C.
 \label{eq:reader-forced-asymmetric}
\end{equation}
The nine-point set and its compact comparison set are
\[
 K_F:=\{P_i^{\rm rad}:0\le i\le5\}\cup\{Q_-,Q_0,Q_+\},
\]
\begin{equation}
 K_{\rm wit}(a,b):=\mathbb B(O,r_*)\cup\{Q_-,Q_0,Q_+\}\subset U_C.
 \label{eq:reader-witness-set}
\end{equation}
These containments follow from \zcref{lem:symmetric-core-witness,lem:asymmetric-core-witness}.

\begin{figure}[!htbp]
\centering
\resizebox{.76\linewidth}{!}{%
\begin{tikzpicture}[scale=3.0,>=Latex]
  \HexCoords
  \DrawHexagon
  \DrawRays

  \foreach \i in {0,...,5}{
    \coordinate (Pr\i) at ($(O)!.16!(V\i)$);
    \fill[DemandRed] (Pr\i) circle (.018);
  }
  \draw[DemandRed!65,line width=.65pt]
    (Pr0)--(Pr1)--(Pr2)--(Pr3)--(Pr4)--(Pr5)--cycle;
  \draw[CenterBlue,dashed,fill=CenterBlue!7,line width=.75pt]
    (O) circle ({.16*sqrt(3)/2});

  \coordinate (Qm) at (-.18,-.50);
  \coordinate (Qz) at (-.02,-.62);
  \coordinate (Qp) at (.17,-.52);
  \fill[GapPurple] (Qm)--++(.025,.043)--++(.025,-.043)--cycle;
  \fill[GapPurple] (Qz)--++(.025,.043)--++(.025,-.043)--cycle;
  \fill[GapPurple] (Qp)--++(.025,.043)--++(.025,-.043)--cycle;
  \node[figlabel,left=2pt] at (Qm) {$Q_-$};
  \node[figlabel,below=2pt] at (Qz) {$Q_0$};
  \node[figlabel,right=2pt] at (Qp) {$Q_+$};
  \node[figlabel,right] at (Pr0) {$P_0^{\rm rad}$};
  \node[figlabel,above right] at (Pr1) {$P_1^{\rm rad}$};
  \node[figlabel,below left] at (Pr4) {$P_4^{\rm rad}$};
  \fill (O) circle (.012) node[above left=1pt,figlabel] {$O$};

  \draw[centerrole]
    (-.49,-.73)--(.50,-.72)--(.02,.24)--cycle;
  \node[figlabel,text=CenterBlue] at (.43,-.63) {putative $U_C$};

  \node[figlabel,atlasbox,align=left,anchor=west] at (1.20,.55)
    {six symmetric points: $P_i^{\rm rad}=(1-c_*)V_i$\\
     three asymmetric points: $Q_-,Q_0,Q_+$\\
     $\mathbb B(0,r_*)\subset\operatorname{conv}\{P_i^{\rm rad}\}$};
  \draw[flow] (1.15,.30)--(.34,.05);

  \node[figlabel,align=center] at (0,-1.22)
    {$K_{\rm wit}=\mathbb B(0,r_*)\cup\{Q_-,Q_0,Q_+\}
      \subset U_C$ under a hypothetical cover};
\end{tikzpicture}%
}
\caption{The zero-gap nine-point witness.  The six symmetric radial points
force a centered disk into the convex C triangle, while the three asymmetric
frontier points supply the missing directional support.  Marker shapes
separate the two constructions: circles are radial witnesses and triangles
are strict-supercritical $AB$-frontier witnesses.  The displayed placement is
illustrative; the point definitions are the exact ones in the preceding
lemmas.}
\label{fig:zero-gap-nine-point-witness}
\end{figure}


\readerheading{The witness set and the comparison set are different.}
The original finite witness is $K_F$, with at most nine points. The set
$K_{\rm wit}$ contains a disk in place of the six radial points and satisfies
\[
 K_{\rm wit}\subseteq\conv(K_F),\qquad
 \Lambda(K_F)=\Lambda(\conv(K_F))\ge\Lambda(K_{\rm wit}).
\]
Thus proving the lower bound for the disk comparison proves it for the
nine-point witness. The disk need not be missed by every V triangle:
its containment in $U_C$ follows from convexity and the six forced radial
points. The later Newton points play the same comparison role; they do not
replace the definitions of the frontier points $Q_-,Q_0,Q_+$.

\readerheading{Active tangent calculation.}
The two outer comparison points use the fixed-start formulas, not the older
junction-start formulas. The exact coefficient comparisons and radius envelope
are specified in \zcref{thm:paper-exact-mixed-certificate}.
The actual disk is unchanged; the smaller scalar radius is used only for
initially checking the tangent residuals.


\subsection{Fixed-start Newton inner reduction}

The three frontier witnesses have already been forced into $U_C$.
We move the outer two toward the junction only for the enclosure
comparison.  These inner points belong to $U_C$ by convexity; they need
not be missed by every V triangle.

Write $U=a+b-1$ and normalize
$\bar\alpha=\alpha/h$, and similarly for the other three coefficients.
For $k=h\lambda$ and $k=h\mu$, the circle restrictions are
\[
 \begin{aligned}
 f_b(k)&=k^2-3(1-U\bar\alpha)k+3b^2-3b+2,\\
 f_a(k)&=k^2-3(1-U\bar\delta)k+3a^2-3a+2.
 \end{aligned}
\]
The preceding root bounds place $1/2$ strictly between the junction
parameter and either first root.  A Newton step from this fixed value gives
\[
 \widehat k_b=\frac{1+3(b-1/2)^2}{2-3U\bar\alpha},\qquad
 \widehat k_a=\frac{1+3(a-1/2)^2}{2-3U\bar\delta}.
\]
The denominators exceed $3/2$.  To keep the line-parameter notation set
$\widehat\xi_-=4\widehat k_b/3$ and
$\widehat\zeta_+=4\widehat k_a/3$.  Define the global points
\[
 \begin{aligned}
 A&=\Psi_4(b-\bar\beta\widehat k_b,
                     \bar\alpha\widehat k_b),\\
 B&=\Psi_4(J),\\
 C&=\Psi_4(\bar\delta\widehat k_a,
                     a-\bar\gamma\widehat k_a).
 \end{aligned}
\]
All subsequent norms and determinants refer to these global vectors.

\begin{lemma}[Fixed-start Newton inner reduction]
\label{lem:technical-newton-reduction}
The local coordinates are rational in $a,b,D$, and the global coordinates
are rational over $\mathbb Q(\sqrt3)$ in those variables.  One has
\[
 A\in(Q_0,Q_-),\qquad C\in(Q_0,Q_+).
\]
Consequently, for the actual disk $\mathcal D_\eta$,
\[
 \widehat K=\conv(\mathcal D_\eta\cup\{A,B,C\})
\]
satisfies
\begin{equation}
 \widehat K\subseteq\conv(K_{\rm wit}(a,b)),\qquad
 \Lambda(K_{\rm wit}(a,b))\ge\Lambda(\widehat K).
 \label{eq:reader-newton-reduction}
\end{equation}
\end{lemma}
\begin{proof}
Each quadratic is positive and decreasing at $k=1/2$, and has positive
second derivative. Its tangent zero is therefore strictly between
$1/2$ and its first root.  The junction is before $1/2$ by
\zcref{lem:fixed-line-signs}.  This proves both open-segment inclusions.
The displayed formulas are rational in the stated fields, and convexity
proves the set inclusion.
\end{proof}

These points lie farther along the two segments than the earlier
junction-start Newton points: for a positive decreasing convex quadratic,
the Newton map has derivative $ff''/(f')^2>0$ before the first root.
The four-contact proof below uses only the segment inclusions and the
unchanged supporting lines; its geometric reasoning is therefore unchanged.
The old algebraic certificate is not reused for the new point formulas.

\subsection{Ordered convex chain and supporting-line bounds}

\begin{proposition}[Four-contact witness geometry]
\label{prop:technical-four-contact-geometry}
Assume \eqref{eq:ab-strict-domain} and $c_*>2/3$. The rays $A,B,C$ occur
in that order in one open $120^\circ$ cone, and
\[
 \operatorname{cross}(B-A,C-B)>0,\qquad
 \|A\|,\|C\|>\eta.
\]
Moreover the outward distances of the lines $AB,BC$ satisfy
\begin{equation}
 d_{AB}:=\frac{\operatorname{cross}(A,B)}{\|B-A\|}\ge h-2\eta,\qquad
 d_{BC}:=\frac{\operatorname{cross}(B,C)}{\|C-B\|}\ge h-2\eta.
 \label{eq:adjacent-overlaps}
\end{equation}
\end{proposition}

\begin{proof}
The local coordinates of all three Newton points lie in $(0,1)^2$.
Centering at $O$ gives coefficient vectors $(u-1,v-1)$ in the two incident
edge directions, so their rays lie in one open $120^\circ$ cone. The two
line identities
\[
 d_{AB}=\alpha(1-b)+\beta>0,\qquad
 d_{BC}=\gamma+\delta(1-a)>0
\]
give their order $A,B,C$. Put
$s_A=h\widehat\xi_- -\lambda_*>0$ and
$s_C=h\widehat\zeta_+ -\mu_*>0$. The frontier directions give
\[
 \operatorname{cross}(B-A,C-B)
 =h s_A s_C(\alpha\delta-\beta\gamma)>0,
\]
since
$\alpha\delta-\beta\gamma=3(1-D)(D+3)/(16\rho)>0$.
For $A=(-X,-Y)$ in the centered cone basis, $X>1/2$, and
\[
 \|A\|^2=(Y-X/2)^2+3X^2/4>3/16.
\]
The reflected calculation gives
\begin{equation}
 \|A\|,\|C\|>h/2>\eta.\label{eq:outer-radii}
\end{equation}
It remains to prove the two supporting-line bounds. Put
$\tau=h-2\eta=h(2c_*-1)$.
We now give the exact analytic verification.  By reflection take $a\le b$ and
set $m=1-b$, $M=1-a$, $U=a+b-1$, $s=m+M$, and $w=M-m$.  The smaller of the
two normalized left sides in \eqref{eq:adjacent-overlaps} is
\begin{equation}
 d=\frac{\mathcal A+U(2m-1)+D(\mathcal A+U)}{2\rho},
 \qquad \mathcal A=1-m+m^2.                      \label{eq:adjacent-d}
\end{equation}
Apply \zcref{lem:bc-capacity-tools} to the complementary pair $(m,M)$.
Its shared deficit bound gives
$2c_*-1\le1-m+m^2=\mathcal A$. On the $C_L$ sheet, put
\[
 X_0=\mathcal A+U,
 \qquad
 Y_0=2\rho\mathcal A-\mathcal A-U(2m-1).
\]
Then
\[
 d-\mathcal A=\frac{D(\mathcal A+U)-Y_0}{2\rho},
\]
and exact expansion gives
\[
\begin{aligned}
 Y_0={}&2(m^2-m+1)U^2+(2m^3-2m+3)U\\
      &+(m^2-m+1)(2m^2-2m+1)>0.
\end{aligned}
\]
Furthermore
\[
 D^2X_0^2-Y_0^2=4m\rho\Phi(U,m),
\]
where
\[
\begin{aligned}
 \Phi(U,m)={}&(1-m)(m^2-m+2)U^2\\
 &+(2-m^2)(m^2-m+1)U\\
 &-m(1-m)^2(m^2-m+1).
\end{aligned}
\]
In these variables
\[
 \chi=U^4-4U^3+5U^2-(m+2)U+m(1-m).
\]
Since $U^2(U^2-4U+5)>0$, the selector $\chi\le0$ implies
$U>m(1-m)/(m+2)$; $\Phi$ is increasing in $U$ and at that lower endpoint is
\[
 \frac{m^2(1-m)(m^4-2m^3+6m^2-6m+4)}{(m+2)^2}>0.
\]
The last quartic equals $(m^2-m)^2+5m^2-6m+4>0$.
Thus $d>\mathcal A\ge2c_*-1$.

On the other radial sheet let $E=\sqrt{4s^2-3}$.  Then
$0\le w<sE$ and $c_*=(s+w)/(1+E)$.  At fixed $U$, differentiation of
\eqref{eq:adjacent-d} can be checked without an implicit estimate.  Namely,
\[
 d=\frac{1+D}{2}-UB(D,w),\qquad
 B(D,w)=\frac{(1+U)(3+D)+w(1-D)}{D^2+3}.
\]
Here $D^2=w^2+6U+3U^2$, so $0\le D'=w/D\le1$, and
\[
 B_w=\frac{1-D}{D^2+3}\ge0.
\]
Writing $N=(1+U)(3+D)+w(1-D)$, direct differentiation gives
\[
 B_D=\frac{(1+U-w)(D^2+3)-2DN}{(D^2+3)^2}>-\frac{34}{27}.
\]
Indeed the first term in the numerator is positive, while
$1+U-w=2a>0$, $N<(7/6)4+1=17/3$, $D<1$, and $D^2+3\ge3$.  If $B_D<0$, then
$B_DD'\ge B_D$; otherwise $B_DD'\ge0$.  Hence
$B'=B_w+B_DD'>-34/27$, and $U<1/6$ yields
\[
 d'<\frac12+\frac16\frac{34}{27}=\frac{115}{162}<1,
\]
whereas
$(2c_*-1)'=2/(1+E)\ge1$.  Hence their difference decreases with $w$.
At the boundary $w=sE$, where $\chi=0$ and $c_*=s$, the preceding sheet
applies: this is an admissible point because $h<s<1$ and
$D^2=4s^4-12s+9<1$.  Indeed
$s^4-3s+2=(s-1)(s^3+s^2+s-2)<0$ on this interval; the second factor is
increasing and already equals $(7h-5)/4>0$ at $s=h$.  The preceding sheet gives
$d\ge2s-1$.  Therefore \eqref{eq:adjacent-overlaps} holds on both
sheets, and reflection supplies the other line.

This proves both supporting-line bounds. Since $0<\eta<h/3$, their
common lower bound $h-2\eta$ is strictly greater than $\eta$, as required
for the exposed-edge geometry used below.
\end{proof}


\subsection{Four supporting contacts and paired radius transfer}
\label{subsec:nine-point-conclusion}

Let $\mathsf R$ and $\mathsf J$ denote counterclockwise rotations through
$2\pi/3$ and $\pi/2$, respectively. For the inner witness define
\[
 H(n)=\max\{\eta,\langle A,n\rangle,\langle B,n\rangle,
                  \langle C,n\rangle\},
 \qquad \Psi(n)=\sum_{j=0}^2H(\mathsf R^jn).
\]
The support formula gives
$\Lambda(\widehat K)=h^{-1}\min_{\|n\|=1}\Psi(n)$.
The minimum can be located without introducing another convex body.

\begin{lemma}[Four-contact enclosure formula]
\label{lem:four-contact-formula}
For the Newton inner witness in the range $c_*>2/3$, put
\[
 n_{AB}=-\frac{\mathsf J(B-A)}{\|B-A\|},\qquad
 n_{BC}=-\frac{\mathsf J(C-B)}{\|C-B\|},
\]
\[
 t_A=\frac{\eta A-\sqrt{\|A\|^2-\eta^2}\,\mathsf JA}{\|A\|^2},
 \qquad
 t_C=\frac{\eta C+\sqrt{\|C\|^2-\eta^2}\,\mathsf JC}{\|C\|^2}.
\]
Then
\begin{equation}
 \Lambda(\widehat K)=\frac1h
 \min\{\Psi(n_{AB}),\Psi(n_{BC}),\Psi(t_A),\Psi(t_C)\}.
 \label{eq:four-contact-minimum}
\end{equation}
\end{lemma}

\begin{proof}
By \zcref{prop:technical-four-contact-geometry}, the rays $A,B,C$ are
ordered in one open $120^\circ$ cone, the chain turns strictly
counterclockwise, and the supporting-line distances exceed $\eta$.
The third point and the disk are strictly on the inward side of each of
$AB,BC$, so both segments are exposed edges. Moreover
\[
 \operatorname{cross}(C-A,B-A)<0,
 \qquad \operatorname{cross}(C-A,-A)>0.
\]
Thus $B$ and the origin are on opposite sides of $AC$, which is not an
exposed edge. The other two straight boundary pieces are the exposed disk
tangents through $A$ and $C$. There is no exposed disk tangent through $B$.
The indicated signs select the tangent before $A$ and the tangent after $C$
in the counterclockwise boundary traversal.

The support-cell rotation lemma \zcref{lem:support-cell-rotation} places
a minimum at these exposed ties or on a disk-only cell. In the latter case
the same value occurs at a boundary tie, since $\|A\|>\eta$.
Rotation through $120^\circ$ leaves just the four displayed normals.
The exposed-hull argument, rather than the general caliper theorem alone,
is what discards the other point and disk ties.
\end{proof}


At the first two contacts the disk supplies the other two supports, so
\begin{equation}
 \Psi(n_{AB})\ge d_{AB}+2\eta\ge h,\qquad
 \Psi(n_{BC})\ge d_{BC}+2\eta\ge h.
 \label{eq:four-contact-line-tests}
\end{equation}
For the two tangencies, put
\[
 \Delta=\operatorname{cross}(C,\mathsf RA)>0,\quad
 \nu=\langle C,\mathsf RA\rangle,\quad N_X=\|X\|^2\quad(X=A,C).
\]
Choosing $C$ as a support source at $\mathsf Rt_A$ and $A$ at
$\mathsf R^2t_C$, with the disk in the other directions, gives
\begin{equation}
 \Psi(t_A)\ge2\eta+
 \frac{\eta\nu+\Delta\sqrt{N_A-\eta^2}}{N_A},\qquad
 \Psi(t_C)\ge2\eta+
 \frac{\eta\nu+\Delta\sqrt{N_C-\eta^2}}{N_C}.
 \label{eq:four-contact-tangent-tests}
\end{equation}
Consequently the only remaining targets are the paired residual signs
\begin{equation}
 P_X(\eta):=(N_X-\eta^2)\Delta^2-
 ((h-2\eta)N_X-\eta\nu)^2\ge0\quad(X=A,C).
 \label{eq:four-contact-residuals}
\end{equation}

% Historical file path retained; the active figure is now the four-contact hull.
\begin{figure}[!htbp]
\centering
\begin{minipage}[t]{.50\linewidth}
\vspace{0pt}
\centering
\resizebox{\linewidth}{!}{%
\begin{tikzpicture}[scale=2.4,>=Latex]
  \node[panellabel,anchor=west] at (-1.02,.72) {(a) the four exposed contacts};
  \fill[CenterBlue!9] (0,0) circle (.19);
  \draw[CenterBlue,line width=.8pt] (0,0) circle (.19);
  \coordinate (A) at (-.70,-.50);
  \coordinate (B) at (0,-.72);
  \coordinate (C) at (.70,-.50);
  \coordinate (L) at (-.14312,.12478);
  \coordinate (T) at (.14312,.12478);
  \draw[GapPurple,line width=1.4pt] (L)--(A)--(B)--(C)--(T);
  \draw[black!35,dashed] (A)--(C);
  \fill (0,0) circle (.012) node[below=2pt,figlabel] {$O$};
  \node[figlabel,text=CenterBlue] at (0,.32) {$\mathcal D_\eta$};
  \foreach \p in {A,B,C}{\fill[GapPurple] (\p) circle (.018);}
  \node[figlabel,below left=2pt] at (A) {$A$};
  \node[figlabel,below=2pt] at (B) {$B$};
  \node[figlabel,below right=2pt] at (C) {$C$};
  \draw[->,DemandRed] (-.35,-.61)--(-.48,-1.03)
    node[below,figlabel] {$n_{AB}$};
  \draw[->,DemandRed] (.35,-.61)--(.48,-1.03)
    node[below,figlabel] {$n_{BC}$};
  \draw[->,RadialTeal] (-.42156,-.18761)--(-.81,.15)
    node[left,figlabel] {$t_A$};
  \draw[->,RadialTeal] (.42156,-.18761)--(.81,.15)
    node[right,figlabel] {$t_C$};
\end{tikzpicture}%
}
\end{minipage}\hfill
\begin{minipage}[t]{.45\linewidth}
\vspace{0pt}
\centering
\textbf{(b) four support sums}\par\medskip
\small
Two consecutive-point contacts:
\[
 \Psi(n_{AB})\ge d_{AB}+2\eta\ge h,
\]
\[
 \Psi(n_{BC})\ge d_{BC}+2\eta\ge h.
\]
Two exposed disk tangencies:
\[
 P_A(\eta),P_C(\eta)\ge0
\]
\[
 \Longrightarrow\quad\Psi(t_A),\Psi(t_C)\ge h.
\]
Hence $\Lambda(\widehat K)\ge1$.
\end{minipage}
\caption{The four-contact terminal for the zero-gap witness. The supporting
lines through $AB$ and $BC$ and the two outer disk tangencies exhaust the
active changes of support source. The dashed segment $AC$ is not exposed:
$B$ and the disk center lie on opposite sides. The drawing is schematic;
the exact witness geometry and the paired tangent residuals are proved in
the text. No auxiliary Minkowski sum or cap covering is needed.}
\label{fig:zero-gap-four-contacts}
\end{figure}

Indeed $P_X(\eta)\ge0$ gives
$\Delta\sqrt{N_X-\eta^2}\ge|(h-2\eta)N_X-\eta\nu|$, which implies the
required lower bound even when the expression inside the absolute value is
negative. No sign is lost by squaring.

The exact certificate uses a smaller rationally specified radius. The
following elementary transfer is the reason those existing signs suffice
for the actual-radius tangencies.

\begin{lemma}[Simultaneous radius transfer]
\label{lem:paired-radius-transfer}
Let $X,Y$ be vectors, $d\ge0$, and $0\le e_0\le1/3$. If
\[
 \|(1-2e_0)X-e_0Y\|\le d,\qquad
 \|(1-2e_0)Y-e_0X\|\le d,
\]
then both inequalities hold at every $e\in[e_0,1/3]$.
\end{lemma}

\begin{proof}
The difference of the two vectors at $e_0$ is $(1-e_0)(X-Y)$, so
\[
 \|X-Y\|\le\frac{2d}{1-e_0}\le3d.
\]
At $e=1/3$ the vectors become $(X-Y)/3$ and $(Y-X)/3$, both in the
closed radius-$d$ ball. Each vector depends affinely on $e$, and the ball
is convex. Interpolation proves both inequalities, including the endpoint
case. This transfers the pair together; it does not assert monotonicity of
either squared residual separately.
\end{proof}

\subsection{Explicit comparisons}
\label{sec:strategy4-exact-certificate}
\let\Fcertificateheading\subsubsection
% Active Case F certificate: fixed-start points and explicit comparisons.
% This source is shared verbatim by the canonical and immediate-proof editions.
The frontier witnesses remain $Q_-,Q_0,Q_+$.  The comparison points $A,B,C$
are the fixed-start Newton points of
\zcref{lem:technical-newton-reduction}; they are not additional forced
frontier witnesses.  The actual disk has radius $\eta=h(1-c_*)$.
Only the tangent-residual calculation uses a smaller scalar radius.
The construction below replaces the former sparse Bernstein transcript,
not the geometric forcing or the two straight-line contact estimates.
It remains an exact computer-assisted argument.
The five new comparison checks use $383$ Chebyshev coefficients, the curvature
comparison uses $113$, and the two increment comparisons use $208$ nonconstant
coefficients; the boundary Taylor and identity checks are additional.
The historical junction-start transcript and its two exact programs remain
unchanged regression material, not inputs for these new point formulas.


\Fcertificateheading{A single rational radius envelope}
Reflect so that $a\ge b$ and set
\[
 \begin{gathered}
 m=1-a,\quad U=a+b-1,\quad a=1-m,\quad b=m+U,\\
 \omega=m(1-m),\qquad b_0=\omega/2.
 \end{gathered}
\]
Here $U$ is a scalar, not an open triangle $U_i$.
The quartic capacity root $C_L(m)$ was defined in
\zcref{eq:core-cstar}.  We first obtain a rational lower bound for
$1-C_L(m)$, then extend it to both capacity branches.

\begin{lemma}[A rational upper-start Newton deficit]
\label{lem:f-rational-newton-deficit}
Let $0<m<1/2$, $h=\sqrt3/2$, and let $C_L(m)$ be the unique root in
$[h,1]$ of $F_m(c)=c^4-c^2+mc-m^2$.
Set $\omega=m(1-m)$ and $c_0=1-\omega/2$. Then
\[
 r(m):=1-c_0+\frac{F_m(c_0)}{F_m'(c_0)}
 =\frac{\omega(4-\omega)(3\omega^2-4\omega+4)}
 {8(4+2m-10\omega+6\omega^2-\omega^3)}
\]
satisfies
\[
0<\omega/2<r(m)\le1-C_L(m)<m.
\]
\end{lemma}
\begin{proof}
One has $c_0\ge7/8>h$ and
\[
 F_m(c_0)=\frac{m^2(1-m)}{16}P(m),\qquad
 P(m)=12-28m+17m^2-11m^3+3m^4-m^5.
\]
On $[0,1/2]$, $P'\le-28+17+3/2=-19/2$ and
$P(1/2)=33/32>0$. Thus $F_m(c_0)>0$.
The polynomial $F_m$ is increasing and strictly convex on $[h,1]$,
so $c_0>C_L(m)$ and its tangent zero
$c_1=c_0-F_m(c_0)/F_m'(c_0)$ satisfies
$C_L(m)\le c_1<c_0$. This proves the first bounds.
If $1-m<h$, then $C_L(m)>1-m$ immediately. Otherwise it follows from
\[
 F_m(1-m)=-m((1-m)^3+m^2)<0.
\]
The displayed rational formula follows by substitution; its denominator
factor $4+2m-10\omega+6\omega^2-\omega^3=2F_m'(c_0)$ is positive.
\end{proof}

\begin{lemma}[Uniform affine radius envelope]
\label{lem:f-uniform-affine-radius-envelope}
\label{lem:paper-branchwise-cbar}
In the strict F domain, put $U=a+b-1$, $m=\min(1-a,1-b)$, and let
$c_*$ be the exact common-pair capacity. For any
$0<r\le1-C_L(m)$, define
\[
 E_r(m,U)=r+(m-r)\left(1-\frac Ur\right)_+
 =\max\left\{r,m-\frac{m-r}{r}U\right\}.
\]
Then
\[
 1-c_*\ge E_r(m,U).
\]
In particular, the rational function in
Lemma~\ref{lem:f-rational-newton-deficit} gives one explicit radius recipe
without testing the original capacity-branch selector.
\end{lemma}
\begin{proof}
Reflect to $a\ge b$ and set $s=1-U$, $C=C_L(m)$. The exact capacity
formula is $c_*=C$ on the quartic branch and
\[
 g_m(s)=\frac{2(s-m)}{1+\sqrt{4s^2-3}}
\]
on the triangular branch. On the latter, its selector
$\chi=F_m(s)>0$ gives $s>h$ because
$\chi\le s^4-3s^2/4$, and hence $s>C$.
Writing $E=\sqrt{4s^2-3}$,
\[
 g_m'(s)=\frac{2(E-3+4ms)}{E(1+E)^2}\le0
 \qquad(C<s\le1).
\]
The root equation gives
$m/C=(1\pm\sqrt{4C^2-3})/2$, and substitution shows $g_m(C)\le C$.
Thus $c_*\le C$ on both branches, so $1-c_*\ge r$.

It remains to consider $0<U<r$. Here $s>1-r\ge C$, so the triangular
formula applies. On the analytic interval $s\in[1-r,1]$,
\[
 g_m''(s)=\frac{8\{s(3+9E-2E^2)-m(3+9E+2E^3)\}}
 {E^3(1+E)^3}>0.
\]
Indeed $m/s\le1/\sqrt3<2/3$ and
\[
 (3+9E-2E^2)-\frac23(3+9E+2E^3)
 =1+3E-2E^2-\frac43E^3\ge\frac23
 \quad(0\le E\le1),
\]
since the last cubic is concave and lies above its endpoint chord.
Endpoint cases follow by continuity. Therefore $U\mapsto g_m(1-U)$
is convex. Its endpoint bounds are
$g_m(1)=1-m$ and $g_m(1-r)\le C\le1-r$.
The chord inequality gives
\[
 c_*\le(1-U/r)(1-m)+(U/r)(1-r),
\]
which is the claimed affine bound. Combining it with $1-c_*\ge r$
proves the result.
\end{proof}


For this choice of $r$ put
\begin{equation}
 \kappa=\frac{m-r}{r},\qquad
 e_0=\max\{r,m-\kappa U\},\qquad e_*=1-c_*.
 \label{eq:f-uniform-radius}
\end{equation}
Thus $0<b_0\le r\le e_0\le e_*$, $r<m$, and $r\le1-h<1/6$.
The slope satisfies
\begin{equation}
 1+\frac m2+2m^2\le\kappa\le3,
 \qquad \ell(m):=\frac{1+5m}{2}\le\kappa.
 \label{eq:f-slope-bounds}
\end{equation}
For the upper bound use $r\ge m(1-m)/2$.  For the lower bound write
$r=N_r/L_r$, using the positive numerator and denominator in
\zcref{lem:f-rational-newton-deficit}.  Exact expansion gives
\[
 mL_r-N_r-(1+m/2+2m^2)N_r=\frac{m^3}{2}H_7(m),
\]
where
\[
 H_7(m)=4+m(164-249m)+m^3(312-238m)
            +m^5(136-45m)+12m^7>0.
\]
Finally $1+m/2+2m^2-\ell(m)=2(m-1/2)^2$.

\Fcertificateheading{One residual and its exact polynomial definition}
Let $X=A$, $Y=\mathsf R^{-1}C$, and set
\[
 N_X=\|X\|^2,\quad N_Y=\|Y\|^2,\quad
 \nu=\langle X,Y\rangle,\quad
 \Delta_0=\frac{\operatorname{cross}(Y,X)}h.
\]
The ordered-chain geometry makes $\Delta_0>0$.  The Gram identity is
\begin{equation}
 h^2\Delta_0^2=N_XN_Y-\nu^2.
 \label{eq:certificate-Gram}
\end{equation}
For $t=1-2e$ define
\begin{equation}
 Q_X(e)=\Delta_0^2-\|tX-eY\|^2,\qquad
 Q_Y(e)=\Delta_0^2-\|tY-eX\|^2.
 \label{eq:certificate-Q}
\end{equation}
These are tangent residuals, not the frontier points $Q_j$.

Use barred coefficients $\bar\alpha=\alpha/h$, and similarly for the
other three frontier coefficients.  Put
\[
 A_1=\rho\bar\alpha,\quad B_1=\rho\bar\beta,\quad
 D_1=\rho\bar\delta,\quad G_1=\rho\bar\gamma,
 \qquad \rho=a^2+ab+b^2,
\]
\[
 n_b=12b^2-12b+7,\quad n_a=12a^2-12a+7,
 \qquad E_b=8\rho-12UA_1,\quad E_a=8\rho-12UD_1,
\]
\[
 \begin{aligned}
 J_b&=(1-2b)E_b+(A_1+2B_1)n_b,& L_b&=E_b-A_1n_b,\\
 J_a&=(1-2a)E_a+(D_1+2G_1)n_a,& L_a&=E_a-D_1n_a.
 \end{aligned}
\]
The denominators satisfy $E_a,E_b>6\rho>0$.
In an orthonormal frame the two vectors are
\[
 X=(-J_b/(2E_b),-hL_b/E_b),\qquad
 Y=(-J_a/(2E_a),hL_a/E_a).
\]
Consequently $4E_a^2E_b^2Q_X(e)$ equals
\[
 \begin{aligned}
 \mathcal N(D,e)={}&(J_bL_a+J_aL_b)^2
 -((1-2e)J_bE_a-eJ_aE_b)^2\\
 &-3((1-2e)L_bE_a+eL_aE_b)^2.
 \end{aligned}
\]
Reducing modulo $D^2=K:=4\rho-3$ defines, by exact polynomial division,
\begin{equation}
 \operatorname{rem}\mathcal N=2\rho^2(\mathcal A+D\mathcal B),\qquad
 \Pi=(1+3K)\mathcal A+K(3+K)\mathcal B.
 \label{eq:f-explicit-polynomials}
\end{equation}
Both $\mathcal B$ and $\Pi$ are quadratic in $e$.  Thus
\[
 Q_X(e)=\frac{\rho^2}{2E_a^2E_b^2}(\mathcal A+D\mathcal B).
\]
No independently stored coefficient table defines these polynomials.

\begin{lemma}[Radius transfer to the actual tangencies]
\label{lem:paper-residual-to-tangency}
If $0<e_0\le e_*<1/3$ and $Q_X(e_0),Q_Y(e_0)\ge0$, then the two
actual-radius tangent contact sums in
\zcref{eq:four-contact-tangent-tests} are at least $h$.
\end{lemma}
\begin{proof}
Apply \zcref{lem:paired-radius-transfer} to $X,Y,\Delta_0$.
This gives $Q_X(e_*),Q_Y(e_*)\ge0$.  The Gram identity gives
\begin{equation}
 P_A(he)=h^2N_XQ_X(e).
 \label{eq:certificate-PA}
\end{equation}
\begin{equation}
 P_C(he)=h^2N_YQ_Y(e).
 \label{eq:certificate-PC}
\end{equation}
The positive determinant sign and the unsquared implication after
\zcref{eq:four-contact-residuals} give the two support bounds.
Only the residual is tested at $e_0$; the disk in the line-contact and
exposed-hull arguments remains the actual disk of radius $he_*$.
\end{proof}

\Fcertificateheading{Explicit low-degree comparisons}
For $s,t\in[-1,1]$ write $m=(1+s)/4$, $U=(1+t)/12$.
The elementary Chebyshev recurrence
\[
 T_0(z)=1,\quad T_1(z)=z,\quad T_{n+1}(z)=2zT_n(z)-T_{n-1}(z)
\]
gives $|T_n(z)|\le1$ on $[-1,1]$.
If $F-g=\sum c_{ij}T_i(s)T_j(t)$, then
\begin{equation}
 |F-g|\le\sum_{i,j}|c_{ij}|.
 \label{eq:f-chebyshev-budget}
\end{equation}
The coefficients are computed over $\mathbb Q$, using the exact power
conversion and then checking reconstruction.  Thus each displayed error
budget below is a finite rational inequality, not a sampling claim.

\begin{lemma}[Explicit algebraic comparison bounds]
\label{lem:f-explicit-comparison-bounds}
On $\mathcal R=[0,1/2]\times[0,1/6]$ the polynomials defined by
\zcref{eq:f-explicit-polynomials} satisfy
\[
 \mathcal B(m,U,b_0)>4,\quad \mathcal B(m,U,m)>210,
 \quad \mathcal B_{ee}<-5600.
\]
Put
\[
 \begin{aligned}
 H_1&=-\Pi_{Uee},&
 H_2&=\Pi_{UUe}(m,U,0)-2\Pi_{Uee},\\
 H_3&=\Pi_{UUe}(m,U,1/3)-2\Pi_{Uee},&
 H_4&=\Pi_{Ue}(m,U,1/3).
 \end{aligned}
\]
Then $H_1>250000$, $H_2>500000$, $H_3>550000$, and $H_4>20000$.
Moreover
\[
 -\Pi_{UU}(m,U,1/3)+2\ell\Pi_{Ue}(m,U,1/3)-\ell^2\Pi_{ee}>21500.
\]
Writing $x=1-2m$, one also has
\[
 \begin{split}
 \Pi_U(m,U,b_0)\ge{}&140+503x+\frac{6265}{2}x^2\\
 &+25000(U-b_0)\qquad(U\ge b_0),
 \end{split}
\]
\[
 \Pi_U(m,U,1/3)>10700\qquad(m\ge1/3).
\]
On the strict physical domain, $N_X\ge N_Y$ and $\Pi_{ee}<0$.
The latter sign also holds on $0<U\le r(m)$.
\end{lemma}
\begin{proof}
We specify the certificate rather than an existence assertion for a sum
of squares.  The four normalized targets $H_i/1000$ have the quadratic
models
\[
 \begin{aligned}
 G_1={}&130s^2+80st-124s+52t^2+100t+454,\\
 G_2={}&466s^2+414st+124s+238t^2+771t+1379,\\
 G_3={}&286s^2+228st-202s+154t^2+348t+1090,\\
 G_4={}&18s^2+4st-14s+6t^2+17t+44.
 \end{aligned}
\]
Their exact Chebyshev error sums are bounded, respectively, by
$70,300,176,8$.  For $G=as^2+bst+cs+dt^2+jt+k$ complete the square as
\[
 G=a\left(s+\frac{bt+c}{2a}\right)^2
       +\gamma(t+1)^2+\delta(t+1)+L,
\]
where
\[
 \gamma=d-\frac{b^2}{4a},\quad
 \delta=j-\frac{bc}{2a}-2\gamma,\quad
 L=k-j+d-\frac{(c-b)^2}{4a}.
\]
For the four models $\gamma,\delta>0$, and
\[
 L=\frac{21188}{65},\quad\frac{373211}{466},\quad
   \frac{210031}{286},\quad\frac{57}{2},
\]
respectively.  Subtracting the error budgets proves the four bounds.

For $\mathcal B(m,U,b_0)$ use the cubic
\[
 \begin{split}
 10g={}&-292s^3+556s^2t+1504s^2+364st^2-317st-1795s\\
       &-28t^3+272t^2-389t+804.
 \end{split}
\]
Its exact error is less than $17/2$.  On $s\le0$, $g_{ss}>0$ and
$g_s(0,t)<0$, so $g(s,t)\ge g(0,t)\ge387/10$.
On $0\le s\le1$, the affine matrix $10(\nabla^2g-8I)$ has positive
upper-left entry and determinant at all four corners.  Its determinant
values are $55687,1026919,82551,23383$.  Convexity of the
positive-semidefinite cone gives $\nabla^2g\succeq8I$.
At $(2/3,1/3)$ the value is $383/30$ and gradient is
$(13/45,-35/18)$.  Completing a square gives
$g\ge1623259/129600>25/2$, hence $\mathcal B(m,U,b_0)>4$.
The other two $\mathcal B$ bounds follow by exact coefficient grouping
in $x=1-2m\in[0,1]$, $6U\in[0,1]$.

For the curvature target divided by $100$, use
\[
 \begin{split}
 g_c={}&76s^3-160s^2t-184s^2+104st^2+226st-132s\\
       &-8t^3-34t^2+283t+788.
 \end{split}
\]
The exact Chebyshev error is less than $55$.
Setting $X_0=1-s$, $Y_0=1+t$ and grouping nonnegative products on
$[0,2]^2$ gives
\[
 g_c\ge277-37X_0+52X_0^2
 =270+52(X_0-37/104)^2+87/208>270.
\]
This proves the curvature bound.  The first-derivative comparison uses
the polynomial increment
\[
 R(m,U)=\frac{\Pi_U(m,U,b_0)-\Pi_U(m,0,b_0)}{U},
\]
with its polynomial value at zero.  Two fixed-chart quadratic lower
models prove $R>25000$; the exact intercept grouping gives
$\Pi_U(m,0,b_0)\ge-2985+503x+12515x^2/2$.
Substituting $b_0=(1-x^2)/8$ proves the stated restricted estimate.
Grouping the $U$ coefficients for $m\ge1/3$ proves the bound $10700$.

For clarity, the two physical-domain signs are not claimed on an
unnecessary larger algebraic domain.  Exact norm subtraction factors as
\[
 N_X-N_Y=\frac{8U(a-b)\rho^2}{E_a^2E_b^2}\mathcal H(D,U).
\]
The explicit power-grouping proof gives
$\mathcal H\ge1261/162$ on $0\le U\le D^2/6$, $0\le D\le1$.
Also the definition of $\Pi$ has the weighted-conjugate identity
\[
 \Pi=\frac{(1+D)^3\mathcal N(D,e)+(1-D)^3\mathcal N(-D,e)}{4\rho^2}.
\]
Each numerator is a constant minus two squares of affine functions of
$e$, so $\Pi_{ee}<0$ for the physical $0<D<1$.
On $0<U\le r$, one has $K(m,U)<1$: at $c=C_L(m)$,
$K(m,1-c)-1=4(c-1)(c^3+c^2+c-2)<0$, and $K$ increases in $U$.

All expansions, power groupings, model reconstruction identities and
rational error bounds just specified are reproduced by the accompanying
\nolinkurl{verify_explicit_stage.py}.  The complete identities and
coefficient bounds, including the two increment models, are printed in
its colocated mathematical supplement.  This verifier uses no optimizer,
no stored SOS witness, no subdivision search, and no floating-point sign
test.  These finite coefficient computations are part of the proof;
they do not audit the preceding geometric forcing lemmas.
\end{proof}

\begin{lemma}[Boundary estimates and coordinatewise monotonicity]
\label{lem:f-boundary-comparisons}
For $0\le m\le1/2$,
\[
 \Pi(m,r,r)>\frac{1063}{1700}.
\]
The affine-branch left endpoints obey
\[
 \Pi(m,0,m)>0\quad(m\le1/3),
\]
\[
 \Pi(m,U,1/3)\ge4096/27
 \quad\bigl(m\ge1/3,\ U\ge(m-1/3)/3\bigr).
\]
Moreover $\Pi_e\ge1152$ on $\mathcal R\times[0,1/3]$.
\end{lemma}
\begin{proof}
At $U=0$, $x=1-2m$, exact factorization gives
\[
 \Pi_e(m,0,1/3)=\frac83(x^2+3)
 (21x^8+130x^6+289x^4+712x^2+144)\ge1152,
\]
and $\Pi_{ee}(m,0,e)<0$.  The positive $\Pi_{Ue}$ bound in
\zcref{lem:f-explicit-comparison-bounds} extends this to the whole
rectangle.  In particular $\Pi$ is increasing in both $U,e$ when
$U\ge b_0$ and $b_0\le e\le1/3$.

Put $q(m)=-9m^2/10+98m/125-3711/100000$.
The exact degree-eight comparison proves $b_0\le q\le r$ for
$1/5\le m\le1/2$.  Thus
\[
 \Pi(m,r,r)\ge
 \begin{cases}
 \Pi(m,b_0,b_0)>103,&m\le1/5,\\
 \Pi(m,q,q)>1063/1700,&m\ge1/5.
 \end{cases}
\]
The second polynomial has degree $27$.  With $x=m-41/100$, its explicit
Taylor coefficient bounds yield
\[
 \Pi(m,q,q)>16/25-10|x|+1700x^2\ge1063/1700.
\]
The first interval and the comparison $q\le r$ also use exact rational
coefficient bounds, reproduced in the same standalone verifier and
printed supplement.  They are not consequences of numerical fitting.
The parameter $q$ is only an auxiliary boundary comparison; it does not
replace $r$, $e_0$, or the actual disk.

For $m\le1/3$, the factor $(2m-1)^2$ can be removed from
$\Pi(m,0,m)$, and grouping higher negative powers under lower positive
powers on $0\le1-3m\le1$ proves positivity.
For $m\ge1/3$, put $x=1-2m\in[0,1/3]$.  Direct expansion gives
$\Pi(m,0,1/3)\ge-96+6688x^2/3$.  Integrating
$\Pi_U(m,U,1/3)>10700$ to $U\ge(1-3x)/18$ gives a decreasing
quadratic whose minimum is $4096/27$.
\end{proof}

\begin{theorem}[Explicit disk-tangent certificate]
\label{thm:paper-exact-mixed-certificate}
On the strict F domain, if $e_0\le1/3$, then
\[
 Q_X(e_0)>0,\qquad Q_Y(e_0)>0.
\]
The assertion uses the fixed-start comparison points, not the old
junction-start Newton points.
\end{theorem}
\begin{proof}
Concavity of $\mathcal B$ and its endpoint bounds give
$\mathcal B(m,U,e)>4$ on $b_0\le e\le m$, hence at $e=e_0$.
The elementary identity
\[
 D(1+3D^2)-D^2(3+D^2)=D(1-D)^3\ge0
\]
therefore gives
\[
 \mathcal A+D\mathcal B\ge\frac{\Pi}{1+3K}.
\]
It suffices to prove $\Pi(m,U,e_0)>0$.

For $U\ge r$, $e_0=r$; the restricted derivative estimate and
$\Pi_{Ue}>0$ give $\Pi_U(m,U,r)>0$.  The boundary estimate at $U=r$
settles this range.
For $U\le r$, put $f(U)=\Pi(m,U,m-\kappa U)$ and
$S_k(e)=\Pi_{UU}-2k\Pi_{Ue}+k^2\Pi_{ee}$.
For $0\le e\le1/3$ and the actual $\kappa\ge1$,
\[
 \partial_eS_\kappa=(1-3e)H_2+3eH_3+2(\kappa-1)H_1>0.
\]
On the integration path, $\Pi_{ee}<0$ and $\Pi_{Ue}>0$, so
$\partial_kS_k<0$ for $k\ge0$.  Consequently
\[
 f''(U)=S_\kappa(e)\le S_\kappa(1/3)
 \le S_{\ell(m)}(1/3)<-21500.
\]
The right endpoint is $U=r$.  The left endpoint is zero if $m\le1/3$;
otherwise it is $U_3=(m-1/3)/\kappa\ge(m-1/3)/3$.
All endpoint values are positive by
\zcref{lem:f-boundary-comparisons}.  Concavity proves positivity
throughout the interval.

Thus $Q_X(e_0)>0$.  Exact subtraction gives
\[
 Q_Y(e)-Q_X(e)=(1-e)(1-3e)(N_X-N_Y)\ge0
 \qquad(0\le e\le1/3),
\]
using the established norm ordering.  This proves both signs.
\end{proof}


% Historical junction-start transcript (not an active input):
% dc46aaf263655d5159ecd3a81db72ee82477951d06172f4743b248df37209485
% BEGIN GENERATED 3105X PROVENANCE
% proof/3XXX_CE0/31XX_Nplus1/310X_all_Vd0/3105X_self_contained_direct_Vd0_nine_point/3105X_computation/mixed_overlap_core_data_00.py 359b219f9deda8586217dcbf9ccc10de0ad48578
% proof/3XXX_CE0/31XX_Nplus1/310X_all_Vd0/3105X_self_contained_direct_Vd0_nine_point/3105X_computation/mixed_overlap_core_data_01.py bfc807f906359636035634f799d74d552cf847fc
% proof/3XXX_CE0/31XX_Nplus1/310X_all_Vd0/3105X_self_contained_direct_Vd0_nine_point/3105X_computation/mixed_overlap_core_data_02.py 92887696ef94aeb93f3d95ab9b16a5388406fff4
% proof/3XXX_CE0/31XX_Nplus1/310X_all_Vd0/3105X_self_contained_direct_Vd0_nine_point/3105X_computation/mixed_overlap_core_data_03.py c2791a8a6add132ad7a17655cbb13987124cb24e
% proof/3XXX_CE0/31XX_Nplus1/310X_all_Vd0/3105X_self_contained_direct_Vd0_nine_point/3105X_computation/mixed_overlap_core_data_04.py 7eeabecbbd17543b728130a1f2bc84f79e11a8d2
% proof/3XXX_CE0/31XX_Nplus1/310X_all_Vd0/3105X_self_contained_direct_Vd0_nine_point/3105X_computation/mixed_overlap_core_data_05.py d676114292e722b60d1b792f615ee88ad58b6982
% proof/3XXX_CE0/31XX_Nplus1/310X_all_Vd0/3105X_self_contained_direct_Vd0_nine_point/3105X_computation/mixed_overlap_core_polynomials.py 07f57b178cf74980f8d956541842d0c4c4b2faa7
% proof/3XXX_CE0/31XX_Nplus1/310X_all_Vd0/3105X_self_contained_direct_Vd0_nine_point/3105X_computation/verify_mixed_overlap_core_derivation.py 0a55211ae37f56e1fc461560b0668a01456c7417
% proof/3XXX_CE0/31XX_Nplus1/310X_all_Vd0/3105X_self_contained_direct_Vd0_nine_point/3105X_computation/verify_global_core_positivity.py 821f9a6dc750fbdb80e02e40282b32cc845c3383
% END GENERATED 3105X PROVENANCE

\begin{theorem}[Witness enclosure]
\label{thm:reader-witness-enclosure}
\label{lem:appendix-witness-enclosure-calculation}

For $0<a,b<1$, $a+b>1$, and $a^2+ab+b^2<1$, the comparison set
$K_{\rm wit}(a,b)$ satisfies
\[
 \Lambda(K_{\rm wit}(a,b))\ge1.
\]
\end{theorem}
\begin{proof}
If $c_*\le2/3$, the centered disk gives
$\Lambda(K_{\rm wit})\ge3(1-c_*)\ge1$. Otherwise
$0<\eta<h/3$, and it suffices to enclose $\widehat K$ by
\zcref{eq:reader-newton-reduction}. The line contacts satisfy
\zcref{eq:four-contact-line-tests}.

\zcref{thm:paper-exact-mixed-certificate} proves the paired residual
signs for the fixed-start points at $0<e_0\le1-c_*$.  The Gram identities
and actual-radius transfer in \zcref{lem:paper-residual-to-tangency}
give $P_A(\eta),P_C(\eta)\ge0$. Thus both tangent sums in
\zcref{eq:four-contact-tangent-tests} are at least $h$ as well.
All four contacts in \zcref{eq:four-contact-minimum} have support sum at
least $h$. Therefore
\[
 \Lambda(K_{\rm wit})\ge\Lambda(\widehat K)\ge1.
\]
\end{proof}

\begin{theorem}[Zero-gap obstruction]
\label{thm:reader-zero-gap-obstruction}
There is no hypothetical cover with
\[
 N_{\rm gap}=0,
 \qquad
 N_+=1.
\]
\end{theorem}

\begin{proof}
The construction gives the compact containment
\(K_{\rm wit}(a,b)\subset U_C\), while
\zcref{thm:reader-witness-enclosure} gives
\(\Lambda(K_{\rm wit}(a,b))\ge1\).  This contradicts
\zcref{lem:compact-cover-margin}.
\end{proof}
